\chapter{Transmission Lines and Standing Wave Ratio}
\label{ch:transmission-lines}

A \textbf{transmission line} (coaxial cable, open-wire ladder line,
twisted pair) carries RF energy from a transmitter to an antenna, or
from an antenna to a receiver, with as little loss and distortion as
possible. Unlike the DC and low-frequency AC circuits of earlier
chapters, a transmission line's physical length is often a significant
fraction of the signal's wavelength, and its behavior must be analyzed
as a distributed, wave-carrying system rather than a simple two-terminal
lumped circuit.

\section{Characteristic Impedance}

Every transmission line has a \textbf{characteristic impedance} $Z_0$,
determined entirely by its physical geometry (conductor spacing,
diameter, dielectric) and, for an ideal lossless line, independent of
its length or the frequency of the signal on it:
\[
Z_0 = \sqrt{\frac{L'}{C'}}
\]
where $L'$ and $C'$ are the line's inductance and capacitance \emph{per
unit length}. Common amateur coaxial cable is manufactured to standard
values of 50\,$\Omega$ (the overwhelming standard for amateur HF/VHF/UHF
work) or 75\,$\Omega$ (common in broadcast and video applications);
open-wire ladder line is typically 300--600\,$\Omega$.

\begin{keyidea}
$Z_0$ is \emph{not} something you measure with an ohmmeter across a
disconnected line's ends (a long line will simply read as an open or
short circuit depending on its length relative to a wavelength) --- it
describes the impedance a wave experiences while traveling \emph{along}
the line, and only becomes directly relevant once the line is driving a
load and carrying a traveling wave.
\end{keyidea}

\section{Matched vs.\ Mismatched Loads}

When a transmission line is terminated in a load whose impedance exactly
equals $Z_0$, all of the power traveling down the line is absorbed by
the load; no energy is reflected. This is a \textbf{matched} condition.
If the load impedance $Z_L$ differs from $Z_0$, some of the incident
power is reflected back toward the source, creating a second wave
traveling in the opposite direction. The interference between the
forward (incident) and reflected waves creates a \textbf{standing
wave} pattern of voltage and current along the line.

\section{Reflection Coefficient and SWR}

The fraction of voltage reflected is the \textbf{reflection
coefficient}, $\Gamma$ (Greek capital gamma), a complex number in
general but often quoted by magnitude:
\[
\boxed{\Gamma = \frac{Z_L - Z_0}{Z_L + Z_0}}
\]
$\Gamma = 0$ for a perfectly matched load; $|\Gamma| = 1$ for a
completely reactive/open/short load (total reflection, no power
absorbed).

The \textbf{standing wave ratio (SWR)}, the ratio of the standing wave
pattern's maximum voltage to its minimum voltage along the line, is
related directly to $|\Gamma|$:
\[
\boxed{\text{SWR} = \frac{V_{\max}}{V_{\min}} = \frac{1+|\Gamma|}{1-|\Gamma|}}
\]
SWR is always $\geq 1$, with 1{:}1 representing a perfect match. Because
it depends only on $|\Gamma|$, SWR (unlike $\Gamma$ itself) carries no
information about the \emph{phase} of the mismatch --- only its
severity.

\begin{examplebox}[SWR from a load mismatch]
A 50\,$\Omega$ line is terminated in a 75\,$\Omega$ resistive load.
Find $\Gamma$ and the SWR.

\emph{Solution.}
\[
\Gamma = \frac{75-50}{75+50} = \frac{25}{125} = 0.2, \qquad
\text{SWR} = \frac{1+0.2}{1-0.2} = \frac{1.2}{0.8} = 1.5.
\]
A 1.5:1 SWR, generally considered a mild, acceptable mismatch.
\end{examplebox}

\section{Power Reflected}

The fraction of incident power that is reflected (rather than
delivered to the load) is $|\Gamma|^2$:
\[
\frac{P_{\text{reflected}}}{P_{\text{forward}}} = |\Gamma|^2.
\]

\begin{examplebox}[Power lost to reflection]
For the 1.5:1 SWR example above ($|\Gamma|=0.2$), find the fraction of
power reflected and the fraction delivered to the load.

\emph{Solution.}
\[
\frac{P_{\text{reflected}}}{P_{\text{forward}}} = 0.2^2 = 0.04\ (4\%), \qquad
\text{delivered} = 96\%.
\]
This illustrates why a modest 1.5:1 SWR is rarely a practical problem:
the actual power lost to reflection is small. Very high SWR (5:1 or
worse) reflects a much larger fraction of power and can additionally
stress transmitter output stages and increase line loss significantly.
\end{examplebox}

\section{Why High SWR Matters}

\begin{itemize}
\item \textbf{Transmitter stress}: many solid-state transmitters
automatically reduce power (``foldback'') or shut down when SWR exceeds
a safe limit, since a severely mismatched load can present the final
amplifier stage with voltage or current levels well outside its safe
operating range.
\item \textbf{Increased line loss}: standing waves increase the peak
current and voltage on the line relative to a matched condition,
increasing $I^2R$ and dielectric losses beyond the line's already
frequency-dependent matched-line loss --- the mismatch loss adds on top
of, and interacts with, the line's intrinsic loss.
\item \textbf{Voltage/current peaks}: at high SWR, the standing wave's
voltage and current maxima can be large enough to cause arcing or
component stress at those points along the line, even if the average
power is modest.
\end{itemize}

\begin{safetynote}
Very high SWR at significant power levels can generate dangerously high
RF voltages at points along a feedline or at an antenna feedpoint,
particularly with unterminated or badly damaged coaxial lines. Address
high SWR readings before operating at higher power levels.
\end{safetynote}

\section{Impedance Along a Mismatched Line}

For a lossless line of length $\ell$ and characteristic impedance
$Z_0$, terminated in load $Z_L$, the impedance seen looking into the
line varies periodically with distance from the load, repeating every
half wavelength:
\[
Z_{\text{in}} = Z_0\,\frac{Z_L + jZ_0\tan(\beta\ell)}{Z_0 + jZ_L\tan(\beta\ell)}, \qquad \beta = \frac{2\pi}{\lambda}.
\]
Two special cases are worth knowing without the full formula:

\begin{itemize}
\item A line exactly \textbf{one-half wavelength} (or any integer
multiple) long repeats the load impedance at its input, regardless of
$Z_0$ --- useful for relocating a matched (or known) impedance to a more
convenient physical location.
\item A line exactly \textbf{one-quarter wavelength} long acts as an
impedance \textbf{inverter/transformer}: $Z_{\text{in}} = Z_0^2/Z_L$.
This \textbf{quarter-wave transformer} principle is widely used to
match a real, resistive antenna feedpoint impedance to a different
feedline impedance using a length of line with an intermediate
characteristic impedance, $Z_0 = \sqrt{Z_{\text{in}}Z_L}$.
\end{itemize}

\section{Line Loss}

Real transmission lines dissipate some power even when perfectly
matched, due to conductor resistance (worse at higher frequency because
of the \textbf{skin effect}, Chapter~\ref{ch:inductors}) and dielectric
loss in the insulating material. Manufacturers publish matched-line
loss in dB per 100\,m (or 100\,ft) at various frequencies; loss
increases with frequency and, for a fixed cable type, with distance.
Mismatch (SWR $>1$) adds additional loss on top of this intrinsic
matched-line loss, an effect that becomes more significant as intrinsic
cable loss increases (which is one reason high SWR is more consequential
on a long run of lossy cable than on a short run of low-loss line).

\begin{quickreview}
\item Characteristic impedance $Z_0 = \sqrt{L'/C'}$ depends on line
geometry, not length or (ideally) frequency.
\item $\Gamma = (Z_L-Z_0)/(Z_L+Z_0)$; SWR $=(1+|\Gamma|)/(1-|\Gamma|)$;
reflected power fraction $=|\Gamma|^2$.
\item High SWR can stress transmitters, increase line loss, and produce
dangerous voltage/current peaks at significant power.
\item A half-wavelength line repeats the load impedance; a
quarter-wavelength line transforms impedance as
$Z_{\text{in}} = Z_0^2/Z_L$.
\item Real lines have intrinsic matched loss (worse at higher
frequency); mismatch adds further loss on top of this.
\end{quickreview}

\begin{practiceproblems}
\item A 50\,$\Omega$ line feeds a purely resistive 100\,$\Omega$ load.
Find $\Gamma$ and SWR.
\item For an SWR of 3:1, find $|\Gamma|$ and the fraction of power
reflected.
\item Find the characteristic impedance of a quarter-wave matching
section needed to match a 50\,$\Omega$ line to a 200\,$\Omega$ antenna
feedpoint.
\item Explain why a transmitter's automatic power foldback at high SWR
is a protective feature rather than a fault.
\item A 1:1 SWR is measured at the transmitter end of a long, lossy
feedline. Does this guarantee the antenna itself is well matched?
Explain, considering line loss.
\end{practiceproblems}
